\clearpage\subsection{Literature Survey}
\subsubsection{Prediction markets and informative participants}
Research on prediction markets suggests that prices may be disproportionately influenced by a relatively small group of informed or disciplined participants. Oliven and Rietz [1], for example, find that active market makers in the Iowa Electronic Markets behave more consistently with rational pricing conditions than typical price takers. This supports the “marginal trader” hypothesis: market accuracy may depend less on the average participant than on the traders willing and able to act against mispricing. Evidence from forecast aggregation is more mixed. Chen et al. [2] find that weighting forecasters according to past performance does not significantly outperform simpler pooling methods in their sports-prediction sample. By contrast, Atanasov et al. [3] show that performance weighting, temporal decay and recalibration can produce forecasts that outperform prediction-market prices, particularly earlier in the life of long-duration questions.

Recent Polymarket research provides more direct evidence of participant heterogeneity. Akey et al. [4] find that approximately 69% of users lose money and that the top 0.1% capture more than half of aggregate gains. Their mixture model classifies a considerably larger group as potentially skilled, while identifying liquidity provision—measured through maker activity—as a strong predictor of profitability. Gómez-Cram et al. [5] employ a more stringent trade-direction randomisation test and classify approximately 3% of accounts as skilled. These accounts exhibit out-of-sample persistence and respond more rapidly to public information. The estimates are not directly comparable because the studies use different samples, null hypotheses and definitions of skill.

Profitability alone, however, is an imperfect measure of informational ability. A profitable wallet may earn spreads, maker incentives or arbitrage returns without contributing superior event information. Evidence from conventional financial markets reinforces the need for careful identification: Barber et al. [6] find that fewer than 1% of Taiwanese day traders earn predictably positive returns after fees, despite substantial apparent in-sample heterogeneity. On-chain analysis must also address multiple wallets controlled by one trader, wash trading and incorrectly inferred trade direction [7], [8].

\subsubsection{From event probability to asset risk}
Event studies compare asset returns around information releases \cite{mackinlay}. Their main difficulty here is separating the event's effect from other news and from expectations already reflected in prices. We will use historical response estimates cautiously and show sensitivity ranges when evidence is weak. An event probability is not, by itself, an equity-price forecast.

Portfolio risk has several meanings. Value at Risk (VaR) is a loss threshold at a chosen confidence level; expected shortfall, also called Conditional Value at Risk (CVaR), describes average loss in the worst tail \cite{rockafellar}. Forecast volatility describes future variation, while realised volatility is measured from subsequent returns \cite{andersen}. These quantities answer different user questions and will be labelled separately.

Option-implied volatility reflects option prices and need not equal the volatility expected under the actual return distribution. The implied--realised variance difference can contain a risk premium \cite{carr}. This is why an apparent variance gap is a research signal, not evidence of an arbitrage opportunity. Forecast comparisons will also account for noisy realised-volatility measurements \cite{patton}.

\subsubsection{AI assistance and the research gap}
Retrieval-augmented generation and tool-using agents provide established patterns for connecting language models to external evidence \cite{lewis,yao}. EventLens will use these patterns to retrieve events and explain model outputs, while deterministic software calculates probabilities and risk.

The project's specific contribution is to test the incremental value of persistent wallet information against anonymous flow, carry the comparison into a bounded equity-risk application, and expose the evidence through an understandable interface. It will connect established methods rather than claim that wallet skill, event studies or AI dashboards are themselves new discoveries.
